% Supplement subsections typed below as "Section~S<n>.<m>"; build.sh checks each number against the built supplement.
% supplement-ref supp:sorting=S3.4
% supplement-ref supp:iia=S1.1
% supplement-ref supp:learned-development=S6.2
% supplement-ref supp:motion=S3.6
% supplement-ref supp:controlled=S3.7
\section{Discussion}\label{sec:discussion}

This paper set out to show that a network can compute with rankings and learn by counting votes. The experiments support that claim and one more. A learning rule that computes no derivative trains ranking layers, and the signal it sends down carries information. Motions of the same kind can also train a real-valued encoder network. But ArrowFlow trails the best tuned classical model on all seventeen small tabular datasets, and training its ranking layers adds only a modest gain. So the contribution is a construction and a learning rule inside permutation space, and the paper makes no accuracy argument.

\paragraph{What the training signal carries.}
The central finding is about the signal. Training helps because of the particular signal that the output layer returns, and matched controls separate that signal from mere movement of the filters (Section~\ref{sec:motion-controls}). If the update were only a useful disturbance, a scramble of the same size would cost nothing. Instead, a scrambled signal is worse than freezing the filters altogether, as a post hoc test across the datasets supports. And on most datasets it leaves an example's nearest neighbors less likely to share its class. Both parts of the signal carry information: which filter a motion reaches, and which way it points. Permuted alignment drew a new random assignment of motions to filters in every batch, and a fixed wrong assignment was not tested (Section~S3.6).

\paragraph{What training adds.}
Training adds a modest gain, and an untrained network already does well. With the same readout, the untrained ArrowFlow comes within one percentage point of ArrowFlow on 12 of the 17 datasets. It is a permutation analog of random-weight models, which fix random hidden units and train only a readout \citep{pao1994learning,huang2006extreme}. It is also an analog of random features for kernel machines \citep{rahimi2007random}. In the cited studies, such models generalized well on the benchmarks tested \citep{huang2006extreme} and matched or outperformed kernel machines there \citep{huang2012extreme,rahimi2007random}. The untrained ArrowFlow differs from them in two ways. Its units are permutations, and its readout is a nearest-neighbor vote rather than a least-squares linear map. Moreover, three of its seven views use the class labels, through a linear discriminant projection.

Read by a strong fixed readout, a Kendall-kernel classifier, the trained hidden rankings are no better than the untrained ones (Section~\ref{sec:controlled}). So training improves the neighborhoods that the nearest-neighbor readout uses. It does not improve the representation that a strong readout sees. Where the gain is large, it can sit in one place. On balance-scale, most of it comes from the rare class of level balances, which the support vector classifier recognizes far more often (Section~\ref{sec:controlled}). This locates the gain. It does not show that the network discovered the physical rule of the balance.

\paragraph{What the trained encoder adds.}
Motions can train more than ranking filters. Turned into target scores, they train the weights of an encoder network across a sort that is never differentiated. As a classifier, the trained encoder is more accurate than the same network untrained on every dataset, significantly on 11 (Section~\ref{sec:learned}). The step is target propagation applied at a discrete interface \citep{bengio2014auto,lee2015difference,friesen2018combinatorial}. It does not relax the sort \citep{grover2019neuralsort,blondel2020fast}, and it does not copy a gradient across it \citep{bengio2013estimating}. Still, with evenly spaced scores the step points the same way as a straight-through estimate for the same targets. The conversion multiplies each coordinate's requested change of position by the gap between neighboring scores. A straight-through estimate that treats position as a linear function of score divides by that gap instead. Other methods also train through an exact discrete step (Section~\ref{sec:related}), among them blackbox differentiation \citep{vlastelica2020differentiation,rolinek2020optimizing} and perturbed optimizers \citep{berthet2020perturbed}. This paper does not compare the conversion with them. Inside ArrowFlow the gain is modest and uneven. One dataset improves significantly and none worsens significantly, but six lose some accuracy. The comparison favors the fixed encoder, because every configuration was tuned for it. The encoder network is one small architecture, and the conversion and its settings were developed on single splits of twelve of the seventeen datasets before these runs (Section~S6.2). The main method therefore keeps the fixed encoder. How well networks that mix ranking layers with ordinary layers can do when both are tuned together is still an open question.

\paragraph{What the Arrow lens delivers.}
The lens does two things. First, it locates context dependence. When the current order counts as a voter, no rule can avoid it. Consider any update rule in the setting of Proposition~\ref{prop:borda-iia} that meets two mild conditions. It never reverses a pair on which the current order and all votes agree, and it neither freezes the filter nor always copies one vote. Every such rule is context dependent. Under the weight condition of Proposition~\ref{prop:borda-iia}, the weighted Borda count that updates a filter violates independence of irrelevant alternatives when the filter's current order counts as a voter. An actual update holds that order fixed. In ArrowFlow's layers it is then not Pareto efficient, so Arrow's theorem does not cover it. A direct construction still shows a violation whenever its votes can move the filter (Section~S1.1). A version of Arrow's theorem without Pareto efficiency, Wilson's theorem, implies the violation once the votes of a batch outweigh the current order \citep{wilson1972pareto}. So a filter's order of two items can depend on a third. The forward pass has no such dependence within a layer.

Second, the voting reading names three stages whose rules can be exchanged, although Arrow's theorem speaks only to the update. The experiments exchanged all three, the view stage only with the prototype readout. The nearest-neighbor readout was more accurate than the prototype readout. Against a footrule median inside ArrowFlow's own layers, the Borda count had the higher mean accuracy on 13 of the 17 datasets, significantly on one. Post hoc, the median learned little. On some datasets it barely moved the filters. On most others it moved them but ended within 0.2 points of untrained filters (Section~S3.7). Across views, with the prototype readout, Borda and majority votes differ by at most $0.6$ percentage points (Section~S4). The lens predicted none of these outcomes. We claim no expressive power from the axiom violation, and the voting-rule study is not a test of Arrow's theorem.

\paragraph{Boundaries.}
A second hidden layer did not improve accuracy at a matched configuration, with the earlier relay or with the corrected one. The rule passes supervision through several ranking layers, but the second layer brought no measured benefit. No experiment showed that the relayed signal improves the first layer. On the seven benchmark datasets, sorting the projected scores made no detectable difference to a nearest-neighbor readout, so there the ranking itself cannot be credited with any gain in accuracy. On the ten further datasets, where the comparison is descriptive, sorting raised the error on six, by up to 2.5 points (Section~S3.4).

\paragraph{Limitations.}
Seventeen small tabular datasets are a narrow basis. The largest training partition has 1,848 rows (Segment). Fifteen test sets give wide intervals, and a difference that is not significant does not show that two models are equal. The candidate grid was small and anchored on default settings. A wider search inside the training data could move Table~\ref{tab:main} in either direction. Settings were chosen by accuracy, and the readout does not weight classes. So a rare class can be neglected (Section~\ref{sec:main-benchmark}). Each study's own limits, such as single fitting seeds and unmatched architectures, are stated with its details in Sections~S3 and~S4.
